% Spis skrótów (opcjonalny) - włącz przełącznikiem \showabbreviationstrue w config/metadata.tex.
% Wpisy układaj alfabetycznie. Tytuł spisu pochodzi z \abbrevname (preamble/polish.tex).
\UnnumberedChapter{\abbrevname}

\begin{abbreviations}
    \item[2FA]   Two-Factor Authentication (uwierzytelnianie dwuskładnikowe)
    \item[AAI]   Authentication and Authorization Infrastructure (infrastruktura uwierzytelniania i autoryzacji)
    \item[ASN]   Autonomous System Number (numer systemu autonomicznego)
    \item[BGP]   Border Gateway Protocol (protokół routingu brzegowego)
    \item[CHF]   Confoederatio Helvetica Franc (frank szwajcarski)
    \item[CT]    Certificate Transparency (publiczny rejestr wystawionych certyfikatów TLS)
    \item[CVE]   Common Vulnerabilities and Exposures (baza znanych podatności bezpieczeństwa)
    \item[D-INFK] Department of Computer Science (Wydział Informatyki ETH Zürich)
    \item[DKIM]  DomainKeys Identified Mail (metoda uwierzytelniania wiadomości e-mail)
    \item[DMARC] Domain-based Message Authentication, Reporting, and Conformance (polityka ochrony domeny pocztowej)
    \item[DNS]   Domain Name System (system nazw domenowych)
    \item[ETH]   Eidgenössische Technische Hochschule Zürich (Politechnika Federalna w Zurychu)
    \item[EXIF]  Exchangeable Image File Format (standard metadanych plików graficznych)
    \item[ISG]   IT Service Group (grupa wsparcia IT wydziału D-INFK)
    \item[OSINT] Open-Source Intelligence (biały wywiad - rozpoznanie ze źródeł jawnych)
    \item[RODO]  Rozporządzenie o Ochronie Danych Osobowych (ang. GDPR)
    \item[SPF]   Sender Policy Framework (mechanizm weryfikacji serwerów nadawczych poczty)
    \item[TLS]   Transport Layer Security (protokół kryptograficzny zabezpieczający komunikację)
    \item[WHOIS] Protokół i bazy rejestracyjne nazw domenowych oraz bloków adresowych IP
    \item[RoE] Rules of Engagement (Reguły Zaangażowania)
\end{abbreviations}
